\section{Note on Sources and Delivery}
For an adult parish assembly celebrating All Saints in the 1962 Roman
Missal on 1 November 2026, with the three Sunday commemoration prayers.
The spoken text ends before this note. It develops the reviewed studies'
interpretation \work{Christ gives rest through humility, purity, and
peace}. Chrysostom's explanation of Christ's yoke and his explicit
connection with the Beatitudes supply the principal argument. Augustine
explains humility, conversion's mourning, mercy, purity and ordered
peace; Chrysostom explains outward reconciliation and the Gospel's
qualifications on persecution. Bede's baptismal robes and acknowledgment
of received gifts inform the account of holiness received. The
connections among these passages, the Mass's petitions and the response
to an insult are the homily's synthesis and application. The example
describes a possible response, not a reported experience or patristic
anecdote. The yoke verses following the Alleluia are identified as
context; they do not enlarge its appointment.

The speech contains 1,236 words in its Pandoc plain-text rendering.
At an estimated unhurried 110--120 words per minute, it takes
10.3--11.2 minutes. It was read silently for sense, paragraph transitions
and oral clarity; no audible rehearsal or timed human delivery is
claimed. Exact sources follow outside the speech. Biblical quotation
uses the historical Douay--Rheims/Challoner English. The commentary
supplies no substitute liturgical translation or newly composed prayer
for recitation. The full study's source records identify the received
texts and their rights; historical source wording is excluded from
the project's license.

\section{References}
\begin{itemize}
\item \work{Missale Romanum}, Vatican editio typica (1962),
All Saints, pp.~702--703: Collect, Lesson, Alleluia, Gospel,
Offertory, Secret, Communion and Postcommunion; the three
commemoration orations, pp.~416--417.
\href{https://media.churchmusicassociation.org/pdf/missale62.pdf}{CMAA facsimile}.
\item \work{Douay--Rheims Bible}, Challoner revision,
\href{https://www.gutenberg.org/ebooks/1581}{Gutenberg 1581}:
Matt.~5:1--12a; 11:28--30; Apoc.~7:2--12; Wis.~3:1--3.
The Offertory is the Missal's selection and adaptation, rather
than a complete consecutive biblical quotation.
\item John Chrysostom, \work{Homilies on Matthew},
XXXVIII.3--5, on Matt.~11:28--30 and the Beatitudes' humility;
XV.7, on peacemaking and persecution. NPNF first series 10,
historical English revised by George Prevost and M.~B. Riddle,
\href{https://www.ccel.org/ccel/schaff/npnf110.html}{CCEL delivery}.
\item Augustine, \work{Our Lord's Sermon on the Mount},
I.1.3; I.2.5--9; I.3.10. William Findlay's historical English
with Schaff's revision and editorial notes, NPNF first series 6,
\href{https://www.ccel.org/ccel/schaff/npnf106.html}{CCEL delivery}.
The editor's notes are distinct from Augustine's argument.
\item Bede, \work{The Explanation of the Apocalypse},
trans.~Edward Marshall (Oxford and London: James Parker, 1878),
book I, on Apoc.~7:9,12, pp.~52--53: baptismal robes,
palms and the acknowledgment of God's gifts.
\end{itemize}
